
\begin{frame}{Qu'est-ce que Spark?}

Un \alert{moteur d'exécution} basé sur des opérateurs de haut niveau, comme Pig.

\vfill

Comprend des opérateurs Map/Reduce, et d'autres opérateurs de second ordre.

\vfill

\alert{Introduit un concept de collection résidente en mémoire} (RDD) 
qui améliore considérablement certains traitements, dont ceux basés sur des itérations.

\vfill
\begin{blocgauche}
De nombreuses librairies pour la fouille de données (MLib), le traitement des graphes,
le traitement de flux (\textit{streaming}).
\end{blocgauche}
\end{frame}

\section{Les RDD}

\begin{frame}{Les \textit{Resilient Distributed Datasets} (RDD)}

C'est le concept central : \alert{Un RDD est une collection (pour en rester à notre vocabulaire) calculée
à partir d'une source de données} (MongoDB, un flux, un autre RDD) .

\vfill

Un RDD peut être marqué comme \alert{persistant}: il est alors placé en mémoire RAM et conservé par Spark.
\vfill

Spark conserve
\alert{l'historique des opérations} qui a permis de constituer un RDD, et la reprise sur panne s'appuie
sur  cet historique afin de reconstituer  le RDD en cas de panne.

\vfill
\begin{blocgauche}
Un RDD est un "bloc" \alert{non modifiable}. Si nécessaire il est \alert{entièrement recaculé}.
\end{blocgauche}
\end{frame}


\begin{frame}{Un \textit{workflow} avec RDD dans Spark}

Des \alert{transformations} (opérateurs comme dans Pig) créent des RDD à partir d'une ou deux sources
de données.

\vfill

\figSlideWithSize{spark-rdd}{10cm}

\vfill
\begin{blocgauche}
Les RDD persistants sont en préservés en mémoire RAM, et peuvent être réutilisés par plusieurs 
traitements.
\end{blocgauche}
\end{frame}

\begin{frame}{Exemple: analyse de fichiers log}
 
On veut analyser le fichier journal (\textit{log}) d'une application dont un des modules (M) est suspect.

\vfill

On construit  un programme qui charge le  \textit{log}, ne conserve
que les messsages produits par le module \textit{M} et les analyse.

\vfill

On peut analyser par produit, par utilisateur, par période, etc.


\vfill

\figSlideWithSize{spark-log}{9cm}

\end{frame}


\begin{frame}[fragile]{Spécification avec Spark}

Première phase pour construire \texttt{logM}

\begin{minted}{scala}
    // Chargement de la collection
    log = load ("app.log") as (...)
    // Filtrage des messages du module M
    logM = filter log with log.message.contains ("M")    
    // On rend logM persistant !
    logM.persist();
\end{minted}

\vfill

Analyse à partir de \texttt{logM}

\begin{minted}{scala}
   // Filtrage par produit
    logProduit = filter logM 
       with log.message.contains ("product P")   
    // .. analyse du contenu de logProduit     
\end{minted}

\end{frame}

\section{Reprise sur panne}

\begin{frame}{Reprise sur panne dans Spark}

Un RDD est une collection \alert{partitionnée}.

\vfill

\figSlideWithSize{spark-failover}{8cm}

\vfill

\begin{blocgauche}
Panne au niveau d'un nœud implique un recalcul basé sur le  fragment $F$  persistant qui précède 
ce nœud dans le workflow.
\end{blocgauche}
\end{frame}

\section{Une session introductive}

\begin{frame}[fragile]{\textit{Dataframes} et \textit{Datasets}}

Un RDD, du point de vue du programmeur, c'est un conteneur d'objets java.

\vfill

Le type précis de ces objets n'est pas connu par Spark. Du coup:

\begin{redbullet}
  \item Tout ce que Spark peut faire, c'est appliquer la sérialisation/désérialisation java 
  \item Aucun accès aux objets grâce à un langage déclaratif n'est possible.
  \item Et donc pas d'optimisation, et la nécessité de tout écrire sous forme de fonctions java. 
\end{redbullet}

\vfill

Depuis la version 1.6: on dispose de RDD améliorés: les \textit{Datasets}. On peut les traiter
comme des tables relationnelles. 

\vfill
\begin{blocgauche}
Finalement, le schéma c'est utile, et le relationnel, c'est bien!
\end{blocgauche}
\end{frame}
